\chapter{Discussion}\label{ch:discussion}

This chapter interprets the evaluation results, situates them within the broader literature, and examines their implications for scientific figure accessibility.
We discuss each research question in turn, beginning with RQ1.

%% ===================================================================
\section{RQ1: What Do the Results Tell Us About VLM Figure Understanding?}\label{sec:disc-rq1}
%% ===================================================================

\subsection{The Proprietary--Open Gap Is Narrowing But Not Closed}

Our leaderboard reveals a 27-point spread from the best model (GPT-5.2, 75.5) to the worst (Gemma-3 4B, 48.5), but the gap between proprietary and open-source leaders is surprisingly small.
Qwen3-VL 235B (70.8) trails GPT-5.2 by only 4.7 points, and the difference between Qwen3-VL 32B (70.4) and Gemini 3.1 Pro (72.7) is not statistically significant at the $p<.05$ level.
This finding contrasts with earlier multimodal benchmarks such as MMMU~\citep{yue2024mmmu}, where proprietary models held a commanding lead over open-source alternatives.
The convergence may reflect rapid improvements in open-source visual instruction tuning, particularly within the Qwen-VL family, which benefits from a dedicated visual encoder trained at scale.

However, the convergence is uneven.
While Qwen models approach proprietary performance, the Gemma-3 family lags 12--27 points behind, and Phi-4 Multimodal underperforms despite its recency.
This suggests that architecture and training data composition matter more than parameter count alone. Qwen3-VL 32B outperforms both Gemma-3 27B (by 11.6 points) and LLaMA-4 Maverick (by 2.6 points), despite similar parameter counts.

\subsection{Numerical Precision as the Fundamental Bottleneck}

The dominance of \emph{Incorrect Numerical Value} errors (30.6\% of all errors) points to a fundamental limitation of current VLMs: they cannot reliably read exact numbers from figure images.
This is not simply an OCR problem. The models often recognise that a value exists but report it imprecisely (e.g., ``approximately 45\%'' when the figure shows 42.3\%).
Our tolerance guidelines already accept $\pm$3 percentage points, meaning the flagged errors exceed even this generous threshold.

This finding aligns with concurrent work on chart question answering.
\citet{masry2022chartqa} report that numerical extraction remains the hardest subtask in ChartQA, and CharXiv~\citep{wang2024charxiv} finds that descriptive questions involving specific values are significantly harder than structural questions.
Our contribution is quantifying this gap in the context of \emph{open-ended description} rather than constrained QA: when models must freely narrate all visible values, numerical errors multiply because each figure contains many readable values.

The implication for accessibility is sobering: a blind researcher relying on VLM-generated descriptions would encounter, on average, 2.7 incorrect numerical values per figure (across both judges).
For scientific figures where exact values carry interpretive significance, this error rate is non-trivial.

\subsection{The Completeness--Accuracy Trade-off}

Models face an implicit trade-off between completeness and accuracy.
A terse description avoids numerical errors but misses visual features; a detailed description covers more atoms but risks more inaccuracies.
Our data suggests that top models resolve this trade-off better: GPT-5.2 achieves both higher atom accuracy (75.0\%) and higher completeness (74.5\%) than Gemma-3 4B (accuracy 55.5\%, completeness 50.3\%).
Weaker models do not fail on one dimension alone; they fail on both.

This contrasts with the assumption that smaller models might achieve reasonable accuracy by being conservative.
In practice, smaller models hallucinate \emph{more} (Gemma-4B: 0.65/fig vs.\ GPT-5.2: 0.14/fig), suggesting that they lack the discriminative capacity to distinguish what is in the image from what their language prior expects to see.

\subsection{The Language Paradox}

German figures score highest (mean 72.5) across all 13 models, surpassing English (62.8).
This counterintuitive result (one might expect English to dominate given its prevalence in VLM training data) has two explanations.

First, \textbf{dataset composition}: German figures in our corpus tend to be simpler, with a mean of 15.8 atoms per figure compared to 19.4 for English.
Under normalised MQM scoring, where penalty is divided by $\text{num\_atoms} \times 5.0$, fewer atoms mean each error costs proportionally more in theory, but in practice the simpler figures also elicit fewer errors.

Second, \textbf{chart type distribution}: the German subset contains more bar charts and fewer pie charts than the English subset.
Since pie charts are substantially harder (mean MQM 45--60 vs.\ 67--76 for bar charts), this compositional difference inflates German scores.

Our cross-lingual controlled comparison (13 parallel figures $\times$ 4 languages) definitively resolves this puzzle.
When the same figure is evaluated in all four languages, the German advantage \emph{reverses entirely}: GPT-5.2 scores 66.2 (EN) vs.\ 55.3 (DE), a 16\% gap.
This confirms that the German advantage in the main leaderboard is a dataset composition artefact (simpler figures with fewer atoms and more bar charts) rather than a genuine language-capability effect.

In the controlled setting, Chinese proves surprisingly resilient (EN--CN gap 4--17\%), likely reflecting strong CJK support in recent VLM training data.
Bulgarian and German both suffer 15--42\% gaps depending on model, with smaller models degrading disproportionately: Gemma-3~4B shows a 42\% EN--DE gap compared to Claude's 2\%.
This suggests that multilingual robustness is a property of scale and alignment rather than architecture alone.

\subsection{On the Reliability of LLM-as-Judge}

Our dual-judge framework produces a key methodological finding: \textbf{per-figure scores disagree substantially (Spearman $\rho \approx 0.6$), but model rankings converge ($\rho \geq 0.95$)}.
This mirrors the ``segment vs.\ system'' distinction well-known in machine translation evaluation~\citep{freitag2021experts}, where segment-level metrics correlate weakly with human judgement but system-level rankings are robust.

The practical implication is that LLM judges should not be trusted for individual figure assessments (a score of 65 on a specific figure carries substantial uncertainty), but they are reliable for comparing models across a corpus of 100+ figures.
Our dual-judge design, averaging two judges with different severity profiles (GPT-4o at 78\% Major, Mistral at 47\% Major), further stabilises the rankings.

The systematic harshness of LLM judges (mean bias $-$12.4 vs.\ humans) is worth noting for future work.
Judges over-penalise high-quality descriptions. GPT-5.2 loses 25.6 points from GPT-4o relative to human scores, while Gemma-3 27B loses only 9.4, suggesting that judges are more sensitive to minor imprecisions in otherwise good descriptions.
This asymmetric bias means that LLM-judge scores compress the true performance range.

\subsection{Pie Charts and the Atom Granularity Problem}

Pie charts score 20--30 points below other chart types.
Our capped-score analysis reveals that ${\sim}11$ points of this gap (out of ${\sim}25$) is a scoring artefact: dense atoms in pie chart descriptions (e.g., ``Breitbart in blue at 14\%, Other in grey at 13\%, Drudge Report in orange at 8\%'') pack multiple verifiable facts into a single sentence, allowing one atom to generate 5--7 errors.
Under our normalisation (penalty / (atoms $\times$ 5.0) $\times$ 100), figures with few atoms and many errors per atom are penalised disproportionately.

The remaining ${\sim}14$ points represent genuine difficulty. Pie charts require VLMs to identify colours, read percentages, match labels to segments, and count slices, all of which are tasks where current models struggle.
This has direct implications for accessibility: pie charts are ubiquitous in scientific publications, yet they receive the least reliable automated descriptions.

Future work could address the granularity problem by splitting dense atoms into finer-grained units (one atom per verifiable fact), though this would increase annotation cost substantially.
